% from ega4-18-fr.tex
\oldpage[IV-4]{155}

étant, il résulte d'abord de (0, 16.1.5) et (0, 16.1.4.1) que la relation $\dim(\sh{O}_{X',x'})\geq 2$ entraîne $\dim(\sh{O}_{X,x})\geq 2$. Inversement, supposons que $\dim(\sh{O}_{X,x})\geq 2$ et que $A$ soit universellement caténaire (5.6.2); alors il en est de même de $\sh{O}_{X,x}$ (5.6.3), donc $\dim(\sh{O}_{X'',x''})=\dim(\sh{O}_{X,x})\geq 2$ en vertu de (5.6.10) et on a donc alors $\dim(\sh{O}_{X',x'})\geq 2$. La même conclusion est valable si $\sh{O}_{X,x}$ est \emph{géométriquement unibranche}, car la fibre du morphisme $X'\to X$ au point $x$ est alors réduite à un seul point, et $\sh{O}_{X',x'}$ est une $\sh{O}_{X,x}$-algèbre \emph{entière}; il suffit donc d'appliquer (0, 16.1.5).
\end{remark}

\begin{lemma}[18.9.7.7]
\label{IV.18.9.7.7}
Soient $A$ un anneau local noethérien, de dimension $\geq 2$, intègre et géométriquement unibranche, $Y=\Spec(A)$, $y$ le point fermé de $Y$. Soient $X$ un préschéma localement noethérien et connexe, $f:X\to Y$ un morphisme plat. Alors, pour toute partie fermée $T\subset f^{-1}(y)$, $X-T$ est connexe.
\end{lemma}

Le même raisonnement qu'au début de (18.9.7.2) ramène à prouver que, pour tout point $x\in f^{-1}(y)$, $\Spec(\sh{O}_{X,x})-\{x\}$ est connexe. On peut donc se borner au cas où $X=\Spec(B)$, $B$ étant un anneau local noethérien et l'homomorphisme $A\to B$ correspondant à $f$ étant local; posons $U=Y-\{y\}$, et notons que $V=f^{-1}(U)$ est dense dans $X$ (2.3.10); il suffit donc de prouver que $V$ est \emph{connexe}. Gardons les notations de (18.9.7.5), et posons donc $A^{(1)}=\Gamma(U,\sh{O}_Y)$, $Y_1=\Spec(A^{(1)})$, $X_1=X\times_Y Y_1=\Spec(B\otimes_A A^{(1)})$, de sorte qu'on a le diagramme commutatif
\[
  \xymatrix{
    X \ar[d]_{f} & X_1 \ar[l]^{g_1} \ar[d]^{f_1} \\
    Y & Y_1 \ar[l]_{g}
  }
\]

En vertu de (2.3.1) et de la définition de $A^{(1)}$, on a un isomorphisme canonique $\Gamma(V,\sh{O}_X)\xrightarrow{\sim}B\otimes_A A^{(1)}$, et il suffit donc de prouver que ce dernier anneau est \emph{local}, un tel anneau ne pouvant être le produit de deux anneaux non réduits à 0 (III, 7.8.6.1). Mais on a vu (18.9.7.5) que $A^{(1)}$ est une sous-$A$-algèbre de la clôture intégrale $A'$ de $A$, puisque $\dim(A)\geq 2$. L'hypothèse que $A$ est géométriquement unibranche entraîne que le morphisme $\Spec(A')\to\Spec(A)$ est radiciel au point $y$ (6.15.3); \emph{a fortiori} $g:Y_1\to Y$ est radiciel au point $y$, et par suite $g_1$ est radiciel au point fermé $x$ de $X$; comme en outre $g_1$ est un morphisme entier, les points fermés de $X_1$ sont au-dessus de $x$, donc il n'existe qu'un seul point fermé de $X_1$, ce qui achève la démonstration de (18.9.7.7).

Il est clair que (18.9.7.7) démontre (18.9.7.3) lorsque $\dim(A)\geq 2$ (même sans supposer les fibres $X_y$ géométriquement réduites).

C) \emph{Cas où $\dim(A)=1$.} On a vu au début de la preuve de (18.9.7.3) que l'on peut supposer $\dim(B\otimes_A k)=0$, et par platitude (6.1.1.1), on a donc
\[
  \dim(B)=\dim(A)+\dim(B\otimes_A k)=1.
\]
Puisque $A$ est intègre, $Y$ se compose de deux points, le point générique $\eta$ et le point fermé $a$. En outre, puisque $f$ est plat, toute composante irréductible de $X$ domine $Y$ (2.3.4), donc $X-\{b\}$, qui est l'ensemble des points maximaux $\xi_i$ de $X$, est égal à

